\Problem{Arbres Verd-Blau}

\Statement

Considerem arbres binaris on cada node pot tenir un color associat:
\emph{verd} o \emph{blau}. Diem que un arbre binari amb colors és un
\emph{arbre verd-blau} quan compleix aquestes dues propietats:

\begin{itemize}

\item Tot node verd té tots els seus fills blaus.

\item Tots els camins que van d'un node fins a qualsevol de les seves fulles té el mateix nombre de nodes blaus.

\end{itemize}


Donat un arbre verd-blau, diem que la seva  \emph{alçada blava} és el nombre de
nodes blaus en els camins des de la seva rel fins a les seves fulles
(com que l'arbre és verd-blau, aquest valor és únic).

\medskip

Per exemple, l'arbre de la figura següent és un arbre verd-blau i la seva alçada blava
és 3.

\FigureC{width=8cm}{arbre}

\Input

L’entrada comença amb un natural $n$, el nombre d’arbres que cal tractar.
Després apareix la descripció dels $n$ arbres, corresponent al seu recorregut en
preordre, amb els nodes verds marcats amb caràcters \texttt{G}, els nodes blaus marcats
amb caràcters \texttt{B} i els arbres buits marcats amb guions.


\Output

Per a cada arbre de l'entrada, si aquest és un arbre verd-blau, cal escriure  la
seva alçada blava. Altrament, cal escriure \texttt{no}.

\Observation

Heu de crear l'arbre on la informació del color sigui explícita en el node, i
vetllar per no tenir fuites de memòria.

\SampleTwoCol
